\subsubsection{LLMListPage}
\label{sec:LLMListPage}
\paragraph{Descrizione}
LLMListPage è la pagina che si occupa di mostrare la lista degli LLM salvati.
\paragraph{SottoComponenti}
LLMListPage è composto dai seguenti sotto-componenti:
\begin{itemize}
    \item \textbf{LLMListComponent}: Il componente con dichiarati i metodi e parametri necessari alla creazione di una LLMListPage tra i quali:
    \begin{itemize}
        \item \textbf{mode}: Parametro settato in fase di routing, da informazioni sull chiamata HTTP che ha chiamato il componente consentendo comportamenti specifici a seconda della chiamata; ciò consente ad un componente di poter gestire leggere variazioni di comportamento conservando la stessa logica di fondo. Tale parametro consente la differenziazione tra la funzione di visualizzare gli LLM e quella di visualizzare gli LLM in fase di test.
        \item \textbf{showAll()}: Metodo che va a chiamare e mostrare la lista di tutti gli LLM salvati.
        \item \textbf{select()}: Consente il reindirizzamento alla pagina dedita all'avvio del test con le informazione sul LLM selezionato.
        \item \textbf{delete()}: Consente di eliminare LLM salvato.
        \item \textbf{addLLM()}: Consente di salvare un LLM. 
        \item \textbf{onDialogResult()}: Metodo per gestire la conferma di azioni che comporteranno modifiche significative.
    \end{itemize}   
    \item \textbf{SenderReqComponent}: Componente che rende sovrascrivibile il metodo sendReq() dedito all'esecuzione di chiamate HTTP.
    \item \textbf{SelectReqComponent}: Componente dotato del metodo sendReq() che eredita da SenderReqComponent, tale componente funge da wrapper per la chiamata HTTP di selezione per la conversione del tipo ritornato dalla chiamata, in un tipo consono alla logica del componente LLMListComponent. E' presente il riferimento in LLMListComponent.
    \item \textbf{DeleteReqComponent}: Componente dotato del metodo sendReq() che eredita da SenderReqComponent, tale componente funge da wrapper per la chiamata HTTP di eliminazione per la conversione del tipo ritornato dalla chiamata, in un tipo consono alla logica del componente LLMListComponent. E' presente il riferimento in LLMListComponent.
    \item \textbf{DialogComponent}: Il componente con dichiarati i metodi necessari alla visualizzazione della finestra di dialog:
    \begin{itemize}
        \item \textbf{dialogStatus}: Metodo booleano che indica la presenza o meno della finestra di dialog.
        \item \textbf{ngOnInit()}: Metodo che si occupa si settare il parametro dialogStatus a False.
        \item \textbf{setStatus()}: Metodo che imposta il valore del parametro dialogStatus.
    \end{itemize}
    \item \textbf{CommandComponent}: Componente che rende sovrascrivibile il metodo execute() dedito all'esecuzione di azioni specifiche.
    \item \textbf{ContinueEventComponent}: Componente dotato del metodo action() che inoltra la richiesta con DeleteReqComponent.
    \item \textbf{CancelEventComponent}: Componente dotato del metodo action() che interrompe l'azione principale.
    \item \textbf{EventsHandlerConfirmationComponent}: Componente dotato del metodo execute() che eredita da CommandComponent, tale componente esegue l'azione del ContinueEventComponent.
    \item \textbf{EventsHandlerAnnullationComponent}: Componente dotato del metodo execute() che eredita da CommandComponent, tale componente esegue l'azione del CancelEventComponent.
\end{itemize}

\paragraph{Tracciamento dei requisiti per LLMListPage}
\begin{tabularx}{\textwidth}{|c|c|}
    \hline
    \textbf{ID} & \textbf{componente} \\
    \hline
        RFF-1.01 & LLMListComponent\\
    \hline
        RFF-1.02 & LLMListComponent\\
    \hline
        RFF-1.03 & SelectReqComponent\\
    \hline
        RFF-1.04 & MessageHandlerComponent\\
    \hline
        RFF-2.01 & LLMListComponent\\
    \hline
        RFF-2.02 & LLMListComponent\\
    \hline
        RFF-2.03 & LLMListComponent\\
    \hline
        RFF-2.04 & DialogComponent\\
    \hline
        RFF-2.05 & DialogComponent\\
    \hline
        RFF-2.06 & LLMListComponent\\
    \hline
        RFF-2.07 & MessageHandlerComponent\\
    \hline
\end{tabularx}